%% kapa: exam=1 round=32 year=2021 period=2 form=A subject=appraisal_law subject_ko=감정평가관계법규 seq=33 number=33
%% source: Q-Net 감정평가사 기출문제 (cst00309) artlSeq=5208397 file=extracted/1008/5208397/A형/2021년도 제32회 감정평가사 1차 2교시 A형.pdf
%% conversion: kordoc 4.24.51 -> tools/gen1.py (auto); review=proofread-sonnet; flags=none
\begin{question}{33}
부동산등기법령상 권리에 관한 등기에 관한 설명으로 옳지 않은 것은?
\begin{choices}
  \item 지방자치단체의 부동산등기용등록번호는 국토교통부장관이 지정·고시한다.
  \item 등기관이 권리의 변경이나 경정의 등기를 할 때에는 등기상 이해관계 있는 제3자의 승낙이 없는 경우에도 부기로 하여야 한다.
  \item 등기관이 전세금반환채권의 일부 양도를 원인으로 한 전세권 일부이전등기를 할 때에는 양도액을 기록한다.
  \item 등기관이 환매특약의 등기를 할 경우 매매비용은 필요적 기록사항이다.
  \item 국가가 등기권리자인 경우, 등기관이 새로운 권리에 관한 등기를 마쳤을 때에는 등기 필정보를 작성하여 등기권리자에게 통지하지 않아도 된다.
\end{choices}
\end{question}
\answer{2}
